\begin{abstract}
% skip one line to make the abstract start with indent

Bug reports submitted through technical support channels may omit reproduction steps, expected behavior, or the context in which a problem occurred. Collecting this information requires clarification between customers and technical teams. Intelligent reporting assistance offers a way to obtain more precise descriptions during report preparation, but its value must be assessed alongside the effort required from the reporter and the accuracy of the resulting information.

This thesis aims to develop and evaluate an intelligent system for software bug reporting that companies can integrate into their service support processes. The proposed system will use a large language model (LLM) to identify missing or ambiguous information and ask service users targeted clarification questions. Through the Model Context Protocol (MCP), it will access documentation describing the service's capabilities and previously resolved problems to inform those questions. The system is designed to interact with users in their chosen language and organize their confirmed answers into a structured report for developers.

The planned evaluation will compare the system with a conventional structured form and with the same assisted form without documentation access. It will examine report completeness and clarity, reporting time, interaction effort, perceived workload, and the preservation of user-confirmed information. The documentation comparison will assess its additional contribution to question relevance and report completeness. The intended contribution is a prototype and empirical evidence about the benefits and limitations of documentation-informed assistance for preparing accurate and useful bug reports.
\end{abstract}
